\section{Introduction}

\begin{quote}
\textit{"It is not the most intellectual of the species that survives; it is not the strongest that survives; but the species that survives is the one that is able best to adapt and adjust to the changing environment in which it finds itself."}
\hspace{5cm} ― Charles Darwin\footnote{This quote is widely attributed to Charles Darwin, but it does not appear verbatim in his writings. The phrasing is believed to originate from Professor Leon C. Megginson, who paraphrased Darwin’s ideas. Despite its frequent misattribution, the quote effectively captures the essence of Darwinian evolution and has since been popularized in both scientific and managerial literature. } 
\end{quote}


Large Language Models (LLMs) have demonstrated remarkable capabilities across a wide range of tasks. Yet, they remain fundamentally static~\citep{luo2025largelanguagemodelagent}, unable to adapt their internal parameters when encountering novel tasks, evolving knowledge domains, or dynamic interaction contexts. As LLMs are increasingly deployed in open-ended, interactive environments, this limitation becomes a critical bottleneck. In such settings, conventional knowledge retrieval mechanisms prove inadequate, giving rise to agents capable of dynamically adapting their perception, reasoning, and actions in real time. This emerging need for dynamic, continual adaptation signals a conceptual shift in artificial intelligence: \textit{from scaling up static models to developing self-evolving agents}.
While established techniques like Supervised Fine-Tuning (SFT) and Reinforcement Learning (RL) provide the mechanisms for improvement, we define self-evolution not merely by the algorithms used, but by the locus of autonomy. Unlike traditional pipelines where human engineers curate data and schedule updates, a self-evolving agent is capable of continuously learning from new data, interactions, and experiences in real-time, leading to systems that are more robust, versatile, and capable of tackling complex, dynamic real-world problems~\citep{tool_tut}.
This shift is currently driving us toward a promising and transformative path to Artificial Super Intelligence (ASI), where the agents not only can learn and evolve from experience with an unpredictable speed but also perform at or above human-level intelligence across a wide array of tasks~\citep{wang2025theoryagentstoolusedecisionmakers}.

Unlike static LLMs, which remain constrained by their inability to adapt to novel and evolving contexts, self-evolving agents are designed to overcome these limitations by continuously learning from real-world feedback.
This progression reshapes our understanding of agents. 
Self-evolving agents, as a core concept, represent a significant step forward in the evolution of intelligent systems, acting as intermediaries that pave the way for more adaptive and autonomous AI, as shown in Figure~\ref{develop}.
Recent research initiatives have increasingly focused on developing adaptive agent architectures capable of continually learning and adapting from experience, such as recent advancements in agent frameworks~\citep{yin2025godelagentselfreferentialagent}, prompting strategies~\citep{fernando2023promptbreeder}, and different optimization ways to evolve. Notwithstanding these advances, existing surveys predominantly address agent evolution as a subsidiary component within comprehensive agent taxonomies. Previous surveys primarily provide systematic overviews of general agent development, while offering limited coverage of self-evolving mechanisms across constrained scenarios in self-evolving agents~\citep {luo2025largelanguagemodelagent, liu2025advanceschallengesfoundationagents}. For example, ~\cite{luo2025largelanguagemodelagent} discuss several ways to evolve, such as self-learning and multi-agent co-evolution, while \cite{liu2025advanceschallengesfoundationagents} explicitly introduce the evolution in terms of different components of agents, such as tools and prompts. Moreover, some studies focus specifically on the evolution of language models themselves \citep{tao2024survey-self-evo-llm}, rather than on the broader concept of agents. However, these works address isolated components rather than the holistic agent system. Therefore, there is no systematic survey devoted
to a dedicated, comprehensive investigation of self-evolving agents as a first-class research paradigm. This
gap has left fundamental questions underexplored: \textbf{\textit{What aspects of an agent should evolve? When should adaptation occur? And how should that evolution be implemented in practice?}}


To the best of our knowledge, this is the first systematic and comprehensive survey focusing on self-evolving agents, offering a clear roadmap for both theoretical inquiry and practical deployment. However, given that this represents a rapidly forming research area where conceptual boundaries are still being actively negotiated within the community, we frame this survey as a guiding synthesis rather than a review of a fully established paradigm. Instead of enforcing rigid boundaries, we aim to structure the heterogeneous mechanisms emerging in the community into a coherent framework.
We organize our analysis around three foundational questions --- \textit{what, when, and how to evolve} --- and provide a structured framework for understanding each. Specifically, we systematically examine individual agent components, including the model, memory, tools and corresponding workflow, investigating their distinct evolutionary mechanisms (what to evolve of agent in Section 3); then we divide existing evolving methods according to different temporal stages with different learning paradigms such as supervised fine-tuning, reinforcement learning and inference-time evolving (when to evolve in Section 4). We finally summarize different signals to guide the evolution of agents, such as textual feedback or scalar rewards, and also different architectures of agents to evolve, such as single-agent and multi-agent evolution 
(how to evolve in Section 5). Furthermore, we review certain evaluation metrics and benchmarks to track existing advancements of self-evolving agents, emphasizing the importance of co-evolution between evaluation and agents (Section 6). We also examine emerging applications in domains such as coding, education, and healthcare, where continual adaptation and evolution are essential (Section 7). Finally, we identify persistent challenges and outline promising research directions to guide the development of self-evolving agents (Section 8). Through this systematic decomposition of self-evolutionary processes across orthogonal dimensions, we provide a structured and practical framework enabling researchers to systematically analyze, compare, and design more robust and adaptive agentic systems. To sum up, our key contributions are as follows:


\begin{itemize}
    \item We establish a unified theoretical framework for characterizing self-evolutionary processes in agent systems, anchored around three fundamental dimensions: what evolves, how it evolves, and when it evolves, providing clear design guidance for future self-evolving agentic systems.

    \item We further investigate the evaluation benchmark or environment tailored for self-evolving agents, highlighting emerging metrics and challenges related to adaptability, robustness, and real-world complexity.

    \item We showcase several key real-world applications across various domains, including autonomous software engineering, personalized education, healthcare, and intelligent virtual assistance, illustrating the practical potential of self-evolving agents.

    \item We identify critical open challenges and promising future research directions, emphasizing aspects like safety, personalization, multi-agent co-evolution, and scalability.


\end{itemize}

In doing so, our survey provides researchers and practitioners with a more structured taxonomy for understanding, comparing, and advancing research of self-evolving agents from different perspectives. As LLM-based agents are increasingly integrated into mission-critical applications, understanding their evolutionary dynamics becomes essential, extending beyond academic research to encompass industrial applications, regulatory considerations, and broader societal implications.

\begin{figure}[t!]
    \centering
    \tikzstyle{my-box}=[
        rectangle,
        rounded corners,
        text opacity=1,
        minimum height=0.1em,
        minimum width=0.1em,
        inner sep=2pt,
        align=left,
        fill opacity=.5,
    ]
    \tikzstyle{leaf}=[my-box]
    \resizebox{\textwidth}{0.72\paperheight}{%
        \begin{forest}
            forked edges,
            for tree={
                grow=east,
                reversed=true,
                anchor=base west,
                parent anchor=east,
                child anchor=west,
                base=left,
                font={\fontsize{5}{5}\selectfont},
                rectangle,
                draw=hidden-draw,
                rounded corners,
                align=left,
                minimum width=0.1em,
                edge+={darkgray, line width=0.8pt},
                s sep=2pt,
                inner xsep=2pt,
                inner ysep=2pt,
                ver/.style={rotate=90, child anchor=north, parent anchor=south, anchor=center},
            },
            where level=1{text width=4em,font={\fontsize{4}{4}\selectfont}}{},
            where level=2{text width=5em,font={\fontsize{4}{4}\selectfont}}{},
            where level=3{text width=4.5em,font={\fontsize{4}{4}\selectfont}}{},
            where level=4{text width=20em,font={\fontsize{4}{4}\selectfont}}{},
            [
            Self-evolving Agent, draw=gray, color=gray!100, fill=gray!15, thick, text=black, ver, [What to evolve, cyan!100, fill=cyan!8,thick, text=black, 
            [Model
            [\parbox{7em}{Policy}[\parbox{36em}{SCA\citep{zhou2025selfchallenginglanguagemodelagents}, Self-Rewarding Self-Improving\citep{simonds2025selfrewardingselfimproving}, SELF\citep{lu2023self}, SCoRe\citep{kumar2024training}, PAG\citep{jiang2025pag}, TextGrad\citep{yellamraju2024textgrad}, AutoRule\citep{wang2025autorule}, SRLM\citep{wang2025selfreasoninglanguagemodelsunfold}}]]
            [\parbox{7em}{Lesson}[\parbox{36em}{SCA\citep{zhou2025selfchallenginglanguagemodelagents}, AgentGen\citep{hu2024agentgen}, Reflexion\citep{shinn2023reflexion}, AdaPlanner\citep{sun2023adaplanner}, SICA\citep{robeyns2025selfimprovingcodingagent}, Self-Refine\citep{madaan2023selfrefine}, Learn-by-interact\citep{su2025learnbyinteractdatacentricframeworkselfadaptive}, RAGEN\citep{wang2025ragen}, DYSTIL\citep{wang2025dystildynamicstrategyinduction}}]]]
            [Context [\parbox{7em}{Memory}[\parbox{36em}{SAGE\citep{liang2024sage}, Mem0\citep{chhikara2025mem0}, MemInsight\citep{salama2025meminsight}, REMEMBER\citep{zhang2024large}, Expel\citep{zhao2024expel}, Agent Workflow Memory\citep{wang2024agent}, Richelieu diplomacy agent\citep{guan2024richelieu}, ICE\citep{qian2024investigate}}]][\parbox{7em}{Prompt}[\parbox{36em}{APE\citep{zhou2022large}, ORPO\citep{yang2023large}, ProTeGi\citep{pryzant2023automatic}, PromptAgent\citep{wang2023promptagent}, REVOLVE\citep{zhang2024revolve}, PromptBreeder\citep{fernando2023promptbreeder}, DSPy\citep{khattab2023dspy}, Trace\citep{wang2023trace}, TextGrad\citep{yellamraju2024textgrad}, SPO\citep{xiang2025self}, LLM-AutoDiff\citep{yin2025llm}, EvoAgent\citep{yuan2024evoagent}}]]]
            [Tool[\parbox{7em}{Creation}[\parbox{32em}{Voyager\citep{wang2023voyager}, Alita\citep{qiu2025alita}, ATLASS\citep{Haque2025ATLASS}, CREATOR\citep{qian2023creator}, SkillWeaver\citep{zheng2025skillweaver}, CRAFT\citep{yuan2024craft}}]][\parbox{7em}{Mastery}[\parbox{36em}{LearnAct\citep{zhao2024learnact}, DRAFT\citep{fromexplorationtomastery}, ToolLLM\citep{qin2023toolllm}, Toolformer\citep{schick2023toolformer}, Gorilla\citep{patil2024gorilla}}]]
            [\parbox{7em}{Selection}[\parbox{36em}{ ToolGen\citep{wang2025toolgen}, AgentSquare\citep{shang2025agentsquare}, Darwin Godel Machine\citep{zhang2025darwin}, COLT\citep{qu2024colt}, TOOLRET\citep{shi2025toolret}, ToolRerank\citep{zheng2024toolrerank}, PTR\citep{gao2024ptr}, SSO\citep{Nottingham2024sso}}]]]
            [Architecture[\parbox{7em}{Single-Agent}[\parbox{29em}{AgentSquare\citep{shang2025agentsquare}, Darwin Godel Machine\citep{zhang2025darwin}, Gödel Agent \citep{yin2025godelagentselfreferentialagent}, AlphaEvolve\citep{novikov2025alphaevolve}, TextGrad\citep{yellamraju2024textgrad}, EvoFlow\citep{zhang2025evoflow}, MASS\citep{zhou2025maasprompt}}]][\parbox{7em}{Multi-Agent}[\parbox{36em}{
                          AFlow\citep{zhang2024aflow},
                          ADAS\citep{hu2024automated}, 
                          AutoFlow\citep{li2024autoflow}, 
                          GPTSwarm\citep{zhuge2024gptswarm}, 
                          ScoreFlow\citep{wang2025scoreflow},
                          FlowReasoner\citep{gao2025flowreasoner},
                          ReMA\citep{wan2025remalearningmetathinkllms},
                          GIGPO\citep{feng2025groupingrouppolicyoptimizationllm}
                        }]
                      ]
                    ]
                ]
                [
                    When to evolve, color=green!80!black!50!, fill=green!5, thick, text=black,
                      [\parbox{8em}{Intra-test-time self-evolution}
                        [\parbox{10em}{ICL}
                          [\parbox{36em}{
                            Reflexion\citep{shinn2023reflexion}, 
                            SELF\citep{madaan2023self_refine}, 
                            AdaPlanner\citep{sun2023adaplanner}, 
                            TrustAgent\citep{hua2024trustagent}
                          }]
                        ]
                        [\parbox{10em}{SFT}
                          [\parbox{36em}{
                            Self-Adaptive LM\citep{zweiger2025self},
                            TTT-NN\citep{hardt2023test},
                            SIFT\citep{hubotter2024efficiently}
                          }]
                        ]
                        [\parbox{10em}{RL}
                          [\parbox{36em}{
                            LADDER\citep{simonds2025ladder},
                            Ttrl\citep{zuo2025ttrl}
                          }]
                        ]
                      ]
                      [\parbox{8em}{Inter-test-time self-evolution}
                        [\parbox{10em}{SFT}
                          [\parbox{36em}{
                            SELF\citep{lu2023self}, 
                            STaR\citep{zelikman2022star}, 
                            Quiet-STaR\citep{zelikman2024quiet}, 
                            SiriuS\citep{zhao2025sirius}, 
                            Chen et al.\citep{chen2025sft}
                          }]
                        ]
                        [\parbox{10em}{RL}
                          [\parbox{36em}{
                            RAGEN\citep{wang2025ragen}, 
                            Learning-Like-Humans\citep{zhang2025learning}, 
                            WebRL\citep{qi2024webrl}, 
                            DigiRL\citep{bai2024digirl}
                          }]
                        ]
                      ]
                ]
                [
                    How to evolve, color=yellow!80!brown, fill=yellow!20!white, thick, text=black
                    [\parbox{8em}{Reward-based Self-Evolution}
                        [Textual Feedback
                            [\parbox{36em}{Reflexion\citep{shinn2023reflexion}, AdaPlanner\citep{sun2023adaplanner}, AgentS2\citep{agashe2025agents2}, SELF\citep{lu2023self}, Self-Refine\citep{madaan2023self_refine}, SCoRe\citep{kumar2024training}, PAG\citep{jiang2025pag}, TextGrad\citep{yellamraju2024textgrad}}]
                        ]
                        [Internal Rewards
                            [\parbox{36em}{CISC\citep{taubenfeld2025cisc}, Self-Ensemble\citep{xu2025selfensemblemitigatingconfidencedistortion}, SRSI\citep{simonds2025selfrewardingselfimproving}, Self-Certainty\citep{kang2025scalablebestofnselectionlarge}, Self-Rewarding Language Models\citep{yuan2025selfrewardinglanguagemodels}}]
                        ]
                        [External Rewards
                            [\parbox{36em}{Self-Train LM\citep{shafayat2025largereasoningmodelsselftrain}, MM-UPT\citep{wei2025unsupervisedposttrainingmultimodalllm}, CoVo\citep{zhang2025consistentpathsleadtruth}, SWE-Dev\citep{du2025swedevevaluatingtrainingautonomous}, SICA\citep{robeyns2025self}, Feedback Friction\citep{jiang2025feedbackfriction}, USEagent\citep{applis2025useagent}, DYSTIL\citep{wang2025dystildynamicstrategyinduction}, OTCPO\citep{wang2025otcpo}, AutoRule\citep{wang2025autorule}, EGSR\citep{zhang2025learning}, LADDER\citep{simonds2025ladder}, RAGEN\citep{wang2025ragen}, SPIRAL\citep{liu2025spiral}}]
                        ]
                        [Implicit Rewards
                            [\parbox{36em}{Reward Is Enough\citep{song2025reward}, Endogenous reward\citep{li2025generalist}}]
                        ]
                    ]
                    [\parbox{10em}{Imitation \& Demonstration Learning}
                        [Self-Generated
                            [\parbox{36em}{STaR\citep{zelikman2022star}, V-STaR\citep{hosseini2024v}, AdaSTaR\citep{koh2025adastar}, STIC\citep{deng2024enhancing}, GENIXER\citep{zhao2024genixer}}]
                        ]
                        [Others
                            [\parbox{30em}{SiriuS\citep{zhao2025sirius}, SOFT\citep{tang2025bridging}, RISE\citep{qu2024recursive}, IoE\citep{li2024confidence}}]
                        ]
                    ]
                    [\parbox{8em}{Population-based \& Evolutionary Methods}
                        [Single Agent
                            [\parbox{36em}{DGM\citep{zhang2025darwin}, GENOME\citep{zhang2025nature}, SPIN\citep{chen2024self}, SPC\citep{chen2025spc}, STL\citep{mendes2025language}}]
                        ]
                        [Multi-Agent
                            [\parbox{36em}{EvoMAC\citep{hu2024self}, Puppeteer\citep{dang2025multi}, MDTeamGPT\citep{chen2025self}, MedAgentSim\citep{almansoori2025self}}]
                        ]
                    ]
                ]
                [
                    Where to Evolve, color=violet!60, fill=violet!15, thick, text=black,
                      [\parbox{12em}{General Domain}
                        [\parbox{10em}{Memory}
                          [\parbox{36em}{
                            Mobile-Agent-E\citep{wang2025mobile}, 
                            MobileSteward\citep{liu2025mobilesteward}
                          }]
                        ]
                        [\parbox{6em}{Curriculum Learning}
                          [\parbox{40em}{
                            WebRL\citep{qi2024webrl},
                            Voyager\citep{wang2023voyager}
                          }]
                        ]
                        [\parbox{6em}{Model-Agent Co-Evolution}
                          [\parbox{36em}{
                            WebEvolver\citep{fang2025webevolver},
                            UI-Genie\citep{xiao2025ui},
                            Absolute-Zero\citep{zhao2025absolute} 
                          }]
                        ]
                      ]
                      [\parbox{12em}{Specialized Domain}
                        [\parbox{6em}{Coding, GUI  Finalicial...}
                          [\parbox{36em}{
                            EvoMac\citep{hu2024self},
                            SICA\citep{robeyns2025self},
                            AgentCoder\citep{dang2025multi},
                            QuantAgent\citep{wang2024quantagent}
                          }]
                        ]
                        [\parbox{20em}{Others}
                          [\parbox{40em}{
                            Arxiv-copilot\citep{lin2024arxivcopilot}, Voyager\citep{wang2023voyager}
                          }]
                        ]
                      ]
                ]
            ]
        \end{forest}
   }
       \caption{Taxonomy of self-evolving agents, in which agents are analyzed along the \textit{what}, \textit{when}, \textit{how}, and \textit{where} dimensions, with selected representative methods and systems annotated at each leaf node.}
    \label{fig:taxonomy}
\end{figure}


\begin{figure}[h]
  \centering
  \includegraphics[width=\linewidth]{Figures/overview.pdf}

\caption{
\textbf{A comprehensive overview of self-evolving agents across key dimensions.}
From left to right and top to bottom, the figure mirrors the organization of Sections~\ref{sec:what}–\ref{sec:evaluation}.
\textbf{What to evolve} (Sec.~\ref{sec:what}) decomposes agent components: model, context, tools, and architecture, showing where evolution operates.
\textbf{When to evolve} (Sec.~\ref{sec:when}) distinguishes intra-test-time versus inter-test-time self-evolution, corresponding to ICL, SFT, and RL paradigms.
\textbf{How to evolve} (Sec.~\ref{sec:how}) summarizes methodological families—reward-based, imitation \& demonstration, and population-based—together with cross-cutting dimensions such as online/offline, on/off-policy, and reward granularity.
\textbf{Where to evolve} (Sec.~\ref{sec:applications}) contrasts general-purpose and domain-specific deployments (e.g., coding, GUI, finance, medical, education).
\textbf{Evaluation} (Sec.~\ref{sec:evaluation}) outlines goals and metrics—adaptivity, generalization, efficiency, safety—and corresponding evaluation paradigms (static, short-horizon, long-horizon).
Overall, the taxonomy maps the survey’s reasoning flow: defining \emph{what, when, and how} to evolve establishes the foundation for evaluating and advancing self-evolving agents.
}

  \label{overview}
\end{figure}

\begin{figure}[h]
  \centering
  \includegraphics[width=\linewidth]{Figures/timeline.pdf}
  \caption{An evolutionary landscape of several representative self-evolving agent frameworks from 2022 to 2025. The figure chronologically organizes major research milestones in the development of self-evolving agents with capabilities such as autonomous planning, tool use, and continual self-improvement. }
  \label{timeline}
\end{figure}

